\section{Présentation du projet}
\label{sec:presentation}

AutoMaze AI vise à reconnaître, à partir d'une image, six types d'usagers de la route : les personnes, les vélos, les voitures, les motos, les bus et les camions. Une fois entraîné, le modèle sera testé sur de nouvelles images. Nous l'utiliserons ensuite dans un simulateur, où la voiture adapte sa conduite selon ce que le modèle a reconnu.

Le sujet porte sur la classification, et elle reste au centre de notre travail. Le CNN et le simulateur sont deux ajouts que nous proposons : ils viennent en plus des livrables du descriptif et n'en remplacent aucun.

Nous laissons de côté la détection d'objets dans une image entière, la vidéo et la physique réaliste. Dans le simulateur, la position des obstacles est déjà connue : le modèle n'a donc qu'à dire de quoi il s'agit.

Au 4 novembre 2026, le simulateur est la partie la plus avancée, alors que le travail sur les données et les modèles commence à peine :

\begin{center}
\begin{tabularx}{\textwidth}{>{\raggedright\arraybackslash}p{4cm}YY}
  \toprule
  \textbf{Partie} & \textbf{État} & \textbf{Reste à faire} \\
  \midrule
  Environnement Colab et Drive & Notebook de vérification exécuté & Vérification par chaque membre, si ce n'est pas fait \\
  Données & Notebook d'exploration à peine commencé & Choix validé, chiffres par classe, recadrages \\
  Modèles et évaluation & Notebooks préparés, aucun résultat & Tout l'entraînement et l'évaluation \\
  Simulateur & Fonctionne en mode simulation, sans vrai modèle & Branchement de notre modèle et comparaison \\
  \bottomrule
\end{tabularx}
\end{center}

Nous sommes trois étudiants, et chacun de nous sera responsable d'un rôle :

\begin{center}
\begin{tabularx}{\textwidth}{>{\raggedright\arraybackslash}p{4cm}>{\raggedright\arraybackslash}p{2.8cm}Y}
  \toprule
  \textbf{Rôle} & \textbf{Responsable} & \textbf{Missions} \\
  \midrule
  Données & \acompleter{Nom} & Exploration, prétraitement, recadrages et découpage \\
  Modèles & \acompleter{Nom} & Classifieurs classiques, CNN, journal des essais \\
  Évaluation et intégration & \acompleter{Nom} & Évaluation finale, analyse des erreurs, export du modèle, lien avec le simulateur \\
  \bottomrule
\end{tabularx}
\end{center}
